% Part: sets-functions-relations
% Chapter: sets
% Section: proofs-about-sets

\documentclass[../../../include/open-logic-section]{subfiles}

\begin{document}

\olfileid{sfr}{set}{prf}
\olsection{সেট সংক্রান্ত প্রমাণ}

\begin{editorial}
এই অংশের পরিবর্তে এখন \verb|content/method/proofs| ব্যবহৃত হয়;
তাই অধ্যায়ের স্বাভাবিক বিন্যাস থেকে এটি বাদ দেওয়া হয়েছে।
\end{editorial}

\begin{explain}
এতক্ষণ যে সেট ও প্রতীকগুলি প্রবর্তন করেছি, সেগুলি সেট নির্দেশ করতে
এবং তাদের মধ্যে সম্পর্ক প্রকাশ করতে সুবিধাজনক সংক্ষিপ্ত রূপ দেয়।
প্রায়ই এই সম্পর্কগুলি সম্বন্ধে দাবি প্রমাণ করাও দরকার হবে। গাণিতিক
প্রমাণের সঙ্গে পরিচিত না হলে এটি তোমার কাছে নতুন হতে পারে। তাই
একটি সহজ উদাহরণ ধাপে ধাপে দেখব। প্রমাণ করব যে, যে-কোনও দুটি
সেট $X$ ও $Y$-এর জন্য সব সময় $X \cap (X \cup Y) = X$ হয়।
সেটের মধ্যে এমন সমতা কীভাবে প্রমাণ করবে? মনে রাখো, সেট শুধু তার
!!{element}s দিয়েই নির্ধারিত হয়; অর্থাৎ, সেটগুলি অভিন্ন হয় যদি এবং
কেবল যদি তাদের !!{element}s একই হয়। তাই এখানে প্রমাণ করতে হবে:
(a) $X \cap (X \cup Y)$-এর প্রতিটি !!{element} $X$-এরও !!{element},
এবং বিপরীত দিকে (b) $X$-এর প্রতিটি !!{element}
$X \cap (X \cup Y)$-এরও !!{element}। অন্য কথায়, আমরা
(a) $X \cap (X \cup Y) \subseteq X$ এবং
(b) $X \subseteq X \cap (X \cup Y)$, উভয়ই দেখাব।

``$X \cap (X \cup Y)$-এর প্রতিটি !!{element}~$z$ $X$-এরও
!!{element}''-এর মতো সাধারণ দাবি প্রমাণ করতে, প্রথমে একটি ইচ্ছামতো
$z \in X \cap (X \cup Y)$ দেওয়া আছে বলে ধরি; তারপর এই অনুমান থেকে
$z \in X$ প্রমাণ করি। এই ধরনটিকে তুমি ``সাধারণ শর্তাধীন প্রমাণ''
নামে জেনে থাকতে পারো। এই প্রমাণে অনুমান ও সিদ্ধান্তে ব্যবহৃত
সংজ্ঞাগুলিও কাজে লাগাতে হবে; যেমন, এখানে ``$\cap$'' ও ``$\cup$''-এর
সংজ্ঞা। তাই (a)-এর যুক্তি হবে এই রকম:

\begin{quote}
(a) প্রথমে দেখাতে চাই $X \cap (X \cup Y) \subseteq X$।
অর্থাৎ, $\subseteq$-এর সংজ্ঞা অনুযায়ী, যে-কোনও~$z$-এর জন্য
$z \in X \cap (X \cup Y)$ হলে $z \in X$ হবে। তাই ধরি,
$z \in X \cap (X \cup Y)$। $z$ দুটি সেটের ছেদের !!{element}
হয় যদি এবং কেবল যদি সে উভয় সেটের !!{element} হয়। তাই পাই,
$z \in X$ এবং $z \in X \cup Y$। বিশেষত, $z \in X$;
এটাই দেখাতে চেয়েছিলাম।
\end{quote}

এতে প্রমাণের প্রথম অর্ধাংশ সম্পূর্ণ হল। লক্ষ করো, শেষ ধাপে এই কথাটি
ব্যবহার করেছি: কোনও অনুমান থেকে একটি সংযোজন ($z \in X$ এবং
$z \in X \cup Y$) পাওয়া গেলে, সেই একই অনুমান থেকে তার প্রতিটি
অংশও পাওয়া যায়। এই নিয়মটি তুমি ``সংযোজন অপসারণ'', অথবা
$\land$Elim নামে জেনে থাকতে পারো। এবার (b) প্রমাণ করি:

\begin{quote}
(b) এবার প্রমাণ করব $X \subseteq X \cap (X \cup Y)$।
অর্থাৎ, $\subseteq$-এর সংজ্ঞা অনুযায়ী, যে-কোনও~$z$-এর জন্য
$z \in X$ হলে $z \in X \cap (X \cup Y)$-ও হবে। ধরি, $z \in X$।
$z \in X \cap (X \cup Y)$ দেখাতে, ``$\cap$''-এর সংজ্ঞা অনুযায়ী
দেখাতে হবে: (i) $z \in X$ এবং (ii) $z \in X \cup Y$।
এখানে (i) আমাদের অনুমানই; তাই তার জন্য আর কিছু প্রমাণ করার নেই।
(ii)-এর জন্য মনে রাখো, $z$ সেটগুলির সংযোগের !!{element} হয় যদি
এবং কেবল যদি সে অন্তত একটি সেটের !!{element} হয়।
$z \in X$ এবং $X \cup Y$ হল $X$ ও $Y$-এর সংযোগ বলে,
এখানে সেই শর্তটি পূরণ হয়। তাই $z \in X \cup Y$। আমরা
(i) $z \in X$ এবং (ii) $z \in X \cup Y$, উভয়ই দেখিয়েছি।
তাই ``$\cap$''-এর সংজ্ঞা থেকে পাই, $z \in X \cap (X \cup Y)$।
\end{quote}

আলোচনাটি কিছুটা দীর্ঘ হল বটে, কিন্তু সেট ও তাদের সম্পর্ক নিয়ে
কীভাবে যুক্তি দিই, তা এতে বোঝা যায়। সাধারণত এত খুঁটিনাটি বলা হয় না;
বিশেষত, সব সংজ্ঞা আবার বলা নাও হতে পারে। আরও উচ্চস্তরের কোনও
পাঠ্যে এই ফলের প্রমাণ অনেক সংক্ষিপ্ত হবে। সেটি এমন দেখতে হতে পারে:
\end{explain}


\begin{prop}[শোষণ]
সমস্ত সেট $X$, $Y$-এর জন্য,
\[
X \cap (X \cup Y) = X
\]
\end{prop}

\begin{proof}
(a) ধরি, $z \in X \cap (X \cup Y)$। তাহলে $z \in X$;
তাই $X \cap (X \cup Y) \subseteq X$।

(b) এবার ধরি, $z \in X$। তাহলে $z \in X \cup Y$-ও হয়,
আর তাই $z \in X \cap (X \cup Y)$-ও হয়।
সুতরাং $X \subseteq X \cap (X \cup Y)$।
\end{proof}

\begin{prob}
$X \cup (X \cap Y) = X$ বিস্তারিতভাবে প্রমাণ করো। তারপর
একটি সংক্ষিপ্ত প্রমাণ দাও। (ইঙ্গিত: $X \cup (X \cap Y) \subseteq X$
দিকটি প্রমাণ করতে বিভিন্ন ক্ষেত্র ধরে প্রমাণ করতে হবে; এই নিয়মের
আরেক নাম $\lor$Elim।)
\end{prob}

\end{document}
